\setcounter{section}{24}
\section{Gobierno Abierto}

Nombre completo: La gobernanza pública y el gobierno abierto. Concepto y principios informadores del gobierno abierto: colaboración, participación, transparencia y rendición de cuentas. Datos abiertos y reutilización. El marco jurídico y los planes de gobierno abierto en España. La Ley 19/2013, de 9 de diciembre, de transparencia, acceso a la información pública y buen gobierno. El Consejo de Transparencia y Buen Gobierno: Real Decreto 919/2014, de 31 de octubre, por el que se aprueba su estatuto. Funciones. La Oficina de Transparencia y Acceso a la Información (OTAI): Funciones. El Portal de Transparencia. Las Unidades de Información y Transparencia (UITS): Funciones. La transparencia y el acceso a la información en las comunidades autónomas y entidades locales.

\localtableofcontents

\subsection{Gobernanza Pública}
    \noindent\fbox{\parbox{\textwidth}{Se denomina Gobernanza pública a la forma de gobernar}}

    Las administraciones públicas deben asegurar un marco normativo estable y adaptado a las necesidades. Tradicionalmente ha sido un sistema burocrático donde se aplican las leyes a favor del ciudadano. Ahora está cambiando hacia un modelo donde el ciudadano es un cliente y no un administrado.

    Supone un cambio de paradigma en las relaciones administrativas. Se propician políticas públicas con participación público-privada.

    La Gobernanza pública se basa en varios conceptos.
    \paragraph{Simplificación administrativa y reducción de cargas}
        La administración pública reduce la necesidad de aportar información al ciudadano para encontrarla automáticamente (autorización mediante). Reduce el coste de funcionamiento. Se reducen obligaciones innecesarias, petición de datos, requerimiento de dociumentos y posibilidad de presentación electrónica de la documentación. De esta forma se  pretende flexibilizar la interacción.

        RD 931/2017 por el que se regula la MEmoria del Análisis de Impacto Normativo

    \paragraph{Inspección de servicios}
        Es una unidad en cada ministerio, normalmente dependiente de la subsecretaría. Se dedica al control interno y evaluación de los servicios de organismos propios y dependientes. Cada ministerio debe tener una carta de servicios, físios y electrónicos. La gobernanza conlleva la inspección y mejora de servicios públicos.

        Debe haber un análisis de las quejas y sugerencias. Es un procedimiento distinto de la reclamación.

        El RD 799/2005 Regula las inspecciones generales de servicios de departamentos ministeriales.

    \paragraph{Atención a la ciudadanía}
        Diseño e impuilso de programas y proyectos para facilitar a la ciudadanía y empresas los servicios públicos a través de canales disponibles. Se normalizan servicios, procedimientos, documentos y el uso de la imagen institucional.

        La web \hyperref{www.administracion.gob.es}{PAGe} es el punto de acceso General, aglutina todas las acciones con el estado.

    \paragraph{Calidad de las administraciones públicas}
        El objetivo en el marco de la gobernanza es favorecer la mejora contínua de la gestión de las administraciones públicas. Se intenta hacer una mejora constante de la calidad.

        Se establece en el RD 951/2005 el marco general para la mejora de la calidad en la Administración General del Estado las siguientes funciones

        \begin{itemize}
            \item Gestión de cartas y servicios
            \item Metodología y guías
            \item Reconocimiento
            \item Informes de evaluación de calidad de servicios
            \item Gestión del conocimiento
        \end{itemize}

    \paragraph{transparencia}
        Es hacer que los gobiernos ofrezcan la información que no les pertenece en exclusiva. El gobierno debe ofrecer a los ciudadanos información sobre el funcionamiento interno. Este modelo de transparencia solo tiene sentido cuando la información es fácilmente accesible.

\subsection{Gobierno Abierto}
    Nace con Obama, es una cultura de gobernanza que promueve principios de transparencia, integridad, rendición de cuentas y participación. Actualmente es el \textbf{V Plan de Gobierno Abierto}, para el periodo 2025 a 2029

    Se basa en una Web 2.0, pero no está ceñido a esta tecnolofía, ni al ámbito de gobierno y administración pública. Se promueve y regula la descarga de datos administrativos.
    % Aquí se te fue la pelota, revisa las clases y el vídeo

    Se considera que la información pública es propiedad de los ciudadanos. Estos deben de poder acceder de manera transparente, fácil y gratuita.

    Con esto se pretenden las siguientes:
    \begin{itemize}
        \item Profundidad en transparencia y rendición de cuentas
        \item Impulsar, fortalecer y mejorar la calidad de la participación
        \item Favorecer la generación de espacios de colaboración
        \item Fortalecer valores éticos y mecanismos de integridad
        \item Sensibilizar a la sociedad y empleados públicos sobre los valores de gobierno abierto
        \item Cumplimiento de Objetivos de Desarrollo Sostenible
        \item Favorecer acciones que favorezcan la inclusión social, la igualdad de accesibilidad universal
    \end{itemize}

    \noindent\fbox{\parbox{\textwidth}{La brecha digital es una contra de los principios de gbierno abierto y de una sociedad decente}}

    \paragraph{Participación}
        Los ciudadanos deben de poder canalizar su opinión y conocimiento para la definición de políticas públicas de su gobierno. El gobierno abierto facilita a los ciudadanos en general y a los expertos y especialistas.

        Los trámites de \textbf{consulta previa} buscan la opinión de ciudadanos y asociaciones ante la elaboración de un proyecto. Mientras que las \textbf{Audiencias e información pública} buscan hacerlo cuando ya se ha iniciado un proyecto.

        \noindent\fbox{\parbox{\textwidth}{El \hyperref{https://transparencia.gob.es/publicidad-activa/por-materias/normativa-otras-disposiciones/plan-anual}{Plan Anual Normativo} incluye todas las normativas que se van a desarrollar por la administracón de turno.}}

    \paragraph{Colaboración}
        La colaboración incluye la participación en la política pública.

    \paragraph{Rendición de cuentas}
        El deber de los servidores públicos es informar, justificar, responsabilizarse púíblica y periodicamente sobre el uso de los fondos asignados y resultados obtenidos. La rendición de cuentas trata mas de la evolución que de la foto.

\subsection{Ley de Transparencia, Acceso a la información y Buen Gobierno}
    \subsubsection{Información pública}
    \subsubsection{Portal de Transparencia}
        \paragraph{La web}
        \paragraph{Límites}
        \paragraph{El Procedimiento de acceso}
        \paragraph{Unidades de información}
        \paragraph{Reclamación de Transparencia y Buen Gobierno}